% ch09-buckets — temporal buckets, discrete-time hazards, and gradient boosting

\section{The horizon target and the bucketed panel}\label{ch09:sec:panel}

Target T3 of \cref{ch:question} asks the question the sourcing workflow
actually poses: will firm $i$ have a deal within the next $\horizon$?
Formally, the object is the fixed-horizon event probability
$p_i(t,\horizon)$ of \cref{eq:ch01:t3}, consumed by binary,
capacity-constrained actions --- enter diligence now, keep on the watch
list, deprioritize. The natural attack is supervised learning on a
bucketed panel: turn the event history into rows ``firm $i$ at time $t$,
did it raise within $\horizon$,'' and hand the rows to a learner that is
good at nonlinearity. This chapter commits to exactly that attack --- with
gradient boosting as the learner --- but insists on doing it as
statistics. The bucketed panel is not a machine-learning convenience
sitting outside the point-process theory of Part II; it \emph{is} a
discrete-time hazard model, with a precise link to the conditional
intensity (\cref{thm:ch09:hazardlink}), a censoring theory
(\cref{ch09:sec:censoring}), a sampling theory
(\cref{ch09:sec:sampling}), and a dependence structure that invalidates
the i.i.d.\ habits that usually travel with the tooling
(\cref{ch09:sec:dependence}). Every labeling shortcut that folklore ML
treats as plumbing --- what to do with windows that cross the end of the
data, whether a firm that died quietly is a negative, whether sliding the
window weekly ``adds data'' --- is a statistical decision with a signed
bias, and this chapter makes each one explicitly.

\begin{definition}[Bucketed panel]\label{def:ch09:panel}
Fix a horizon $\horizon > 0$ and an evaluation grid of bucket start times
\begin{equation}\label{eq:ch09:grid}
\mathcal{T}_{\horizon} \;=\; \{\, t_b = t_0 + b\,\horizon \;:\; b = 0, 1,
\dots, B_{\horizon}-1 \,\} \subset [0, \Tend],
\end{equation}
aligned to the weekly calendar of \cref{ch:data}. The panel for horizon
$\horizon$ is the row set
$\mathcal{P}_{\horizon} = \{ (i, t_b) : i \in \firms,\; \atrisk_i(t_b) = 1 \}$
--- one row per tracked firm at risk at each bucket start --- with binary
label
\begin{equation}\label{eq:ch09:label}
y_{i,b} \;=\; \ind{\, N_i(t_b + \horizon) - N_i(t_b) \,\ge\, 1 \,},
\end{equation}
and feature vector $\zfeat_{i,b}$ assembled at $t_b$ under
\cref{prot:ch02:pit}:
\begin{enumerate}[label=(\roman*)]
\item the LOCF firm covariates $\tilde\xcov_i(t_b)$ of
\cref{eq:ch02:locf};
\item the staleness age $\staleage_i(t_b) = t_b - \tau_i(t_b)$ --- ``it
has been $\staleage$ since the last deal'' --- carried \emph{always},
as mandated by design decision D8.1 of \cref{ch:stale};
\item engineered history: deal counts $N_i(t_b) - N_i(t_b - w)$ over
trailing windows $w \in \{26, 52, 104\}$ weeks, last round type and size
(with an undisclosed flag), raised-to-date;
\item sector activity summaries: trailing sector deal counts and their
changes, computed from the \emph{complete} sector event stream (including
out-of-universe deals, whose counts are trustworthy even when their firms
are not);
\item as-of macro $\xcov^{\mathrm{mac}}_{\mathrm{pub}}(t_b)$, joined on
publication date.
\end{enumerate}
Negative rows are all at-risk firm-buckets with no deal in the window;
there is no other source of negatives. The committed horizon set is
$\horizon \in \{26, 52, 104\}$ weeks.
\end{definition}

The label \cref{eq:ch09:label} deserves one immediate remark: it depends
on the event history only through a bucket \emph{count}, never through a
within-bucket position. This is deliberate. Defect (C4) of \cref{ch:data}
says interval counts are trusted while positions are not; the panel is
built on exactly the trusted functional, and \cref{cor:ch09:defect} below
prices what that choice costs.

\paragraph{Defect declaration (per D2.1).} Against the ledger of
\cref{tab:ch02:defects}: \textbf{(C1)} corrected --- buckets whose window
crosses $\Tend$ are flagged censored and either dropped or retained with
the exposure-offset likelihood of \cref{ch09:sec:censoring}; they are
never labeled negative. \textbf{(C2)} robust by construction --- no row
exists before $\entry_i$; the deeper selection (the tracked universe is
conditioned on having raised at least once) restricts the estimand to
tracked firms and is owned by \cref{ch:universe}. \textbf{(C3)}
vulnerable with signed bias --- unlabeled deaths sit in the negative
class; mitigations and the residual downward bias are stated in
\cref{ch09:sec:censoring}. \textbf{(C4)} robust by design --- labels are
invariant to within-bucket date jitter (only events within jitter
distance of a bucket boundary can flip a label), at the price of zero
information about within-bucket timing, a price this chapter chooses
deliberately (\cref{cor:ch09:defect}). \textbf{(C5)} vulnerable,
mitigated --- LOCF covariates enter only alongside their staleness age,
and the tree learner can represent the interaction ``informativeness of
$\tilde\xcov$ decays with $\staleage$'' that \cref{ch:stale} shows is the
first-order correction; residual attenuation remains and is inherited
from \cref{ch:stale}. \textbf{(C6)} robust for the stated estimand ---
labels are own-firm events of tracked firms, and out-of-universe deals
enter only through complete sector counts; the panel makes no claim about
untracked firms.

\begin{decision}{D9.1 --- the bucketed panel is a hazard dataset}
The T3 panel is built exactly as \cref{def:ch09:panel}: calendar-aligned
disjoint buckets per horizon, labels from bucket counts only, features
under \cref{prot:ch02:pit} with the staleness age mandatory, censored
buckets flagged rather than zero-labeled, and no rows before first
observation. Rationale: each clause is the discrete-time image of a
continuous-time censoring or truncation rule
(\cref{thm:ch09:hazardlink,prop:ch09:censorbias}); violating any one of
them injects a bias whose sign we can already state, which is worse than
a variance problem because more data will not fix it. Reversal trigger:
none for the labeling rules; the feature list may grow freely under
\cref{prot:ch02:pit}.
\end{decision}

\begin{decision}{D9.2 --- horizon set $\{$6M, 1Y, 2Y$\}$}
We commit to $\horizon \in \{26, 52, 104\}$ weeks and build one model per
horizon. Rationale: the three consuming actions named in
\cref{ch:question} (diligence now; pipeline within a year; deprioritize
if nothing plausible in two years) fix the horizons --- the target is
chosen by the action, not by the data; and per-horizon models keep the
cloglog consistency of \cref{thm:ch09:hazardlink} exact.
Reversal trigger: if the decision layer of \cref{ch:decision} starts
consuming horizon \emph{curves} $\horizon \mapsto p_i(t,\horizon)$
rather than three scalars, the per-$\horizon$ design yields to the
single-model AFT alternative of \cref{ch09:sec:aft}.
\end{decision}

\section{The bucket model is a hazard model}\label{ch09:sec:hazardlink}

We now prove the claim that organizes the chapter: the panel of
\cref{def:ch09:panel} is a discrete-time hazard model, and the bucket
probability is an exact functional of the same conditional intensity that
Part II estimates in continuous time. Nothing about this is asymptotic;
the identity holds bucket by bucket.

\begin{lemma}[Waiting-time hazard of the bucket]\label{lem:ch09:waiting}
Fix a bucket $(i,b)$ with $\atrisk_i(t_b)=1$ and let
$W_{i,b} = \inf\{u > 0 : N_i(t_b+u) > N_i(t_b)\}$ be the waiting time to
firm $i$'s next deal. Suppose the conditional law of $W_{i,b}$ given
$\F_{t_b}$ is absolutely continuous on $(0,\infty)$ with hazard function
$r_{i,b}(u)$. Then
\begin{equation}\label{eq:ch09:hazardmass}
p_{i,b} \;\defeq\; \Prob\big( y_{i,b} = 1 \st \F_{t_b} \big)
\;=\; 1 - \exp\big( -\Lam_{i,b} \big),
\qquad
\Lam_{i,b} \;\defeq\; \int_0^{\horizon} r_{i,b}(u) \dd u ,
\end{equation}
and the hazard is the predictable projection of the conditional
intensity onto the no-event paths:
$r_{i,b}(u) = \E\big[ \lam_i(t_b+u) \st \F_{t_b},\, W_{i,b} > u \big]$.
\end{lemma}

\begin{proof}
The first display is the hazard representation of an absolutely
continuous survival law, $\Prob(W > \horizon \st \F_{t_b}) =
\exp(-\int_0^{\horizon} r)$, applied to
$\{y_{i,b}=0\} = \{W_{i,b} > \horizon\}$. For the second, the innovation
theorem \citep[II.4.2]{abgk1993} states that the intensity of $N_i$ with
respect to the coarser filtration generated by $\F_{t_b}$ and the
post-$t_b$ own-event history is the conditional expectation of
$\lam_i$ given that filtration; on the event $\{W_{i,b} > u\}$ the
post-$t_b$ own-event history is trivial, and the hazard of the first
point of a counting process equals its intensity before the first point.
\end{proof}

The lemma already says something practical: the bucket probability
depends on the future only through an \emph{integrated} intensity, with
the unobserved evolution of covariates and competing events averaged out
along the no-event paths. What the panel model estimates is therefore not
an exotic new object; it is the compensator increment of Part II, coarsened.
To get the classical grouped-data form we add the proportional-hazards
structure.

\begin{assumption}[Grouped proportional hazards within
buckets]\label{asm:ch09:groupedph}
For every bucket $(i,b)$ there is a baseline $\lam_{0,b}(u) \ge 0$,
common across firms, and a coefficient vector $\bcoef$ such that
$r_{i,b}(u) = \lam_{0,b}(u)\, e^{\bcoef\T \zfeat_{i,b}}$ for
$u \in (0,\horizon]$: covariates are frozen at their bucket-start values,
and their proportional effect on the no-event intensity holds throughout
the bucket. Piecewise-constant intensity ($\lam_{0,b}(u)$ constant in
$u$) is the special case used for exposure offsets below.
\end{assumption}

\begin{theorem}[Hazard link]\label{thm:ch09:hazardlink}
Let the panel be built as in \cref{def:ch09:panel} and let
$h_{i,b} = p_{i,b}$ denote the discrete-time hazard
\cref{eq:ch09:hazardmass}.
\begin{enumerate}[label=(\roman*)]
\item \emph{(Exact cloglog / grouped PH.)} Under
\cref{asm:ch09:groupedph},
\begin{equation}\label{eq:ch09:cloglog}
\Prob\big(y_{i,b}=1 \st \zfeat_{i,b}\big)
= 1 - \exp\!\big(-\Lam_{0,b}\, e^{\bcoef\T \zfeat_{i,b}}\big),
\qquad
\operatorname{cloglog} p_{i,b}
\;\defeq\; \log\big(\!-\!\log(1-p_{i,b})\big)
= \alpha_b + \bcoef\T \zfeat_{i,b},
\end{equation}
with $\alpha_b = \log \Lam_{0,b}$, $\Lam_{0,b} = \int_0^{\horizon}
\lam_{0,b}$, and $\bcoef$ \emph{identical} to the continuous-time
coefficients \citep{prenticegloeckler1978,allison1982}. The Bernoulli
likelihood of the panel is the exact likelihood of the point process
observed through the coarsening ``at least one event per bucket'': a
bucket classifier estimates a deterministic coarsening
$\lam \mapsto (\Lam_{i,b})_b$ of the same intensity object as Part II.
\item \emph{(Continuum limit.)} Drop \cref{asm:ch09:groupedph}. If
$u \mapsto r_{i,b}(u)$ is right-continuous at $0$ then
$h_{i,b}/\horizon \to r_{i,b}(0^+) = \E[\lam_i(t_b) \st \F_{t_b}]$ as
$\horizon \downarrow 0$; and for a fixed window $[t, t+v]$ partitioned
into buckets of mesh $\horizon \to 0$, with hazards refreshed at each
bucket start and bounded above,
\[
\prod_{b} \big(1 - h_{i,b}\big) \;\longrightarrow\;
\exp\Big( -\!\int_t^{t+v} \lam_i(u) \dd u \Big)
\quad \text{along no-event paths:}
\]
the product-limit of the discrete hazards recovers the continuous-time
survival factor, so the discrete model converges to the continuous
intensity model, not merely to an approximation of it.
\item \emph{(Logit versus cloglog.)} Writing $\Lam = \Lam_{i,b}$ for the
bucket's hazard mass,
\begin{equation}\label{eq:ch09:linkgap}
\operatorname{logit} p \;=\; \log\big(e^{\Lam} - 1\big)
\;=\; \operatorname{cloglog} p \;+\; d(\Lam),
\qquad
d(\Lam) = \log\frac{e^{\Lam}-1}{\Lam}
= \frac{\Lam}{2} + \frac{\Lam^2}{24} - \frac{\Lam^4}{2880} + \cdots ,
\end{equation}
so for small per-bucket probabilities the log-odds equal the
log-hazard-mass up to $O(\Lam)$: a logistic-link panel model estimates
the same object as the cloglog model to first order. The divergence is
material at large $p$: $d = 0.026$ nats at $p=0.05$, $0.114$ at
$p=0.2$, $0.367$ at $p=0.5$.
\end{enumerate}
\end{theorem}

\begin{proof}
(i) Substituting \cref{asm:ch09:groupedph} into
\cref{eq:ch09:hazardmass} gives $\Lam_{i,b} = e^{\bcoef\T\zfeat_{i,b}}
\Lam_{0,b}$ and hence both displays of \cref{eq:ch09:cloglog}, since
$\operatorname{cloglog}(1 - e^{-\Lam}) = \log \Lam$. For the likelihood
claim: the coarsened data for bucket $(i,b)$ is the binary $y_{i,b}$,
whose conditional law given $\F_{t_b}$ is Bernoulli$(p_{i,b})$ by
\cref{lem:ch09:waiting}; over disjoint buckets the joint likelihood
factorizes by the tower property into exactly the product of these
Bernoulli terms, which is the person-period likelihood of
\citet{allison1982}. No approximation enters at any step.
(ii) $h = 1 - \exp(-\int_0^{\horizon} r)$ with $\int_0^{\horizon} r =
\horizon\, r(0^+) + o(\horizon)$ by right-continuity, and
$1 - e^{-x} = x + O(x^2)$, give $h/\horizon \to r(0^+)$; the
identification of $r(0^+)$ is \cref{lem:ch09:waiting}. For the product
limit, $\log \prod_b (1-h_b) = \sum_b \log(1-h_b) = -\sum_b h_b +
O\big(\sum_b h_b^2\big)$; along a no-event path each refreshed hazard
mass satisfies $h_b = \int_{t_b}^{t_b+\horizon} \lam_i(u)\dd u +
o(\horizon)$, so $\sum_b h_b \to \int_t^{t+v} \lam_i$, while
$\sum_b h_b^2 \le (\max_b h_b) \sum_b h_b \to 0$ because
$\max_b h_b = O(\horizon)$ under the boundedness assumption.
(iii) $p = 1 - e^{-\Lam}$ gives $p/(1-p) = e^{\Lam}-1$, hence the first
equality. For the expansion, $(e^{\Lam}-1)/\Lam = e^{\Lam/2}\,
\sinh(\Lam/2)/(\Lam/2)$, so $d(\Lam) = \Lam/2 + \log\big(\sinh y / y\big)$
with $y = \Lam/2$, and $\log(\sinh y/y) = y^2/6 - y^4/180 + O(y^6)$
yields the stated series. The numerical values follow from
$\Lam = -\log(1-p)$.
\end{proof}

Part (iii) has a sharp practical reading. When per-bucket probabilities
are small --- short horizons, ordinary firms --- the logistic link is the
cloglog link plus a distortion of order the probability itself, and a
logistic learner (including XGBoost with its native
\code{binary:logistic} objective) estimates a hazard-consistent object.
When per-bucket probabilities are large --- the two-year horizon on
recently active firms --- the gap $d(\Lam)$ is covariate-dependent, so it
is \emph{not} absorbed by an intercept: logistic coefficients cease to be
transferable across horizons, while cloglog coefficients remain the same
$\bcoef$ at every $\horizon$ by \cref{eq:ch09:cloglog}. This is the
precise sense in which cloglog, not logit, is the aggregation-coherent
link, and it is why the calibration layer of \cref{ch09:sec:calibration}
is mandatory rather than cosmetic: it repairs, on the validation era,
exactly the link distortion we can bound here.

\begin{corollary}[Choosing $\horizon$ is choosing the aggregation
defect]\label{cor:ch09:defect}
Any two conditional intensities with equal bucket masses
$(\Lam_{i,b})_{i,b}$ induce the same panel law; consequently the panel
Fisher information is identically zero in every parameter direction that
moves $\lam$ only within buckets. In particular, timescale parameters
faster than $\horizon$ (a kernel decay $\decay$ with half-life
$\ll \horizon$) are unidentifiable from the panel.
\end{corollary}

\begin{proof}
The panel likelihood in \cref{thm:ch09:hazardlink}(i) depends on the
underlying process only through $(\Lam_{i,b})$; a parameter direction
that leaves all bucket masses fixed leaves the likelihood fixed, so the
score in that direction vanishes identically and so does its variance.
\end{proof}

This is the $\branch$--$\decay$ ridge of \cref{ch:censoring} returned as
a design dial. There, weekly binning was an observation defect (C4) that
destroyed the decay while sparing the branching scale --- bucketing at or
below one kernel half-life was safe, four half-lives destroyed timing
(\cref{app:results}). Here the same mathematics appears with the sign of
agency reversed: T3 does not \emph{want} within-horizon timing, so we
choose $\horizon$ to integrate it out, and \cref{cor:ch09:defect} states
exactly what has been discarded. The defect is not incurred by
negligence; it is purchased, at a known price, because the consuming
action is horizon-bound.

\section{Censoring and eligibility}\label{ch09:sec:censoring}

Three eligibility rules determine which $(i, t_b)$ pairs become rows and
what their labels may say. Each is the discrete image of a
continuous-time censoring statement, and each has a signed bias when
violated.

\paragraph{Right censoring at the window end (C1).} Define the exposure
fraction of bucket $(i,b)$ as
$e_{i,b} = \min\{1, (\Tend - t_b)/\horizon\}$ (rows require
$t_b < \Tend$, so $e_{i,b} \in (0,1]$). A bucket with $e_{i,b} < 1$ ---
any window crossing June 2024 --- is \emph{censored}: the label
\cref{eq:ch09:label} is not computable from data, because the window's
unobserved remainder may contain the event.

\begin{warning}\label{warn:ch09:censorzero}
Labeling censored buckets $y=0$ is not a conservative convention; it is a
biased estimator with a known sign. It converts ``unobserved'' into
``negative'' and deflates every fitted hazard, most strongly in the most
recent calendar eras --- which are exactly the eras closest to deployment
and the eras whose macro covariates differ most from the training bulk,
so the bias masquerades as an era effect and contaminates the macro
coefficients. This is the bucketed instance of
\cref{warn:ch02:openspells}.
\end{warning}

\begin{proposition}[Censored-bucket bias]\label{prop:ch09:censorbias}
Suppose bucket $(i,b)$ has hazard mass $\Lam_{i,b} > 0$ distributed with
constant within-bucket intensity, and exposure $e = e_{i,b} < 1$. The
naive label $\tilde y_{i,b} = \ind{\text{event observed in }
[t_b, \Tend)}$ satisfies
$\E[\tilde y_{i,b} \st \F_{t_b}] = 1 - e^{-e \Lam_{i,b}}
< 1 - e^{-\Lam_{i,b}} = p_{i,b}$:
on the cloglog scale the fitted linear predictor is shifted by
$\log e < 0$, i.e.\ hazards are biased down by exactly the unobserved
exposure fraction.
\end{proposition}

\begin{proof}
With constant within-bucket intensity the observed sub-window carries
hazard mass $e\Lam_{i,b}$, so the observed-event probability is
$1 - e^{-e\Lam_{i,b}}$ by \cref{lem:ch09:waiting} applied to the
sub-window; $\operatorname{cloglog}(1 - e^{-e\Lam}) = \log \Lam + \log e$.
\end{proof}

Two legal treatments follow. \emph{Drop:} for the pure classification
pipeline, delete rows with $e_{i,b} < 1$. This restores unbiasedness ---
administrative censoring at fixed $\Tend$ is independent by construction
(\cref{ch:data}, defect C1) --- at the cost of discarding the most recent
eras; the cost grows with the horizon, and at $\horizon = 104$ weeks it
amounts to the final two years of bucket starts. \emph{Keep with an
exposure offset:} retain censored rows with the survival-style Bernoulli
likelihood implied by \cref{prop:ch09:censorbias},
\begin{equation}\label{eq:ch09:offsetlik}
\loglik \;=\; \sum_{(i,b)} \Big[\, \tilde y_{i,b} \log\Big(1 -
\exp\big(-e_{i,b}\, e^{\eta_{i,b}}\big)\Big) \;-\; \big(1 -
\tilde y_{i,b}\big)\, e_{i,b}\, e^{\eta_{i,b}} \,\Big],
\end{equation}
where $\eta_{i,b}$ is the cloglog linear predictor (or boosted margin)
and $\log e_{i,b}$ enters as a fixed offset. Under within-bucket constant
intensity this is exact; it is the grouped-data analogue of keeping open
spells with their survival factor. The classification stack commits to
the drop rule for simplicity, and \cref{ch09:sec:aft} names the offset
likelihood and AFT as the routes that recover the discarded eras.

\paragraph{Entry (C2).} No row is emitted for $t_b < \entry_i$: a firm
contributes person-periods only after its first observation. This is
delayed entry handled mechanically by the at-risk indicator, exactly as
in \cref{eq:ch02:jacod}; emitting earlier rows would populate the panel
with firms observed \emph{because} they later raised, a truncation bias
with the opposite sign to \cref{warn:ch09:censorzero}.

\paragraph{Unlabeled exits (C3).} The negative class is contaminated:
a firm that quietly died before $t_b$ has $\atrisk_i(t_b) = 1$ in the
naive risk set and contributes an eternal string of structural negatives.
The consequences and their measured severity are inherited from
\cref{ch:data} and \cref{app:results}: hazards deflate by the
contamination fraction, and covariates correlated with staleness acquire
\emph{inverted} weight, because the Bayes-optimal response to a pool
containing zombies is to bet against staleness. Three mitigations are
committed, in order of force. First, \emph{exposure windows}: a row
$(i,t_b)$ requires observed activity (a deal, or a point-in-time exit
proxy refresh) within the trailing $E$ years. Because the rule is a
predictable function of the observed history, it is exactly independent
censoring; its benefit is not repairing the mechanism but capping each
dead firm's phantom exposure at $E$ years. Second, an explicit
\emph{active model} --- a point-in-time classifier of ``still alive'' ---
used as a feature, not as a hard filter: the rehearsal showed a
discriminative exit model (AUC $0.877$) is \emph{insufficient} to
reconstruct truthful risk sets, so it may inform but must not silently
redefine the estimand. Third, the staleness age itself partially absorbs
the defect: dead firms accumulate enormous $\staleage$, and a flexible
learner routes their predicted probability toward zero through the
$\staleage$ feature --- which is simultaneously why the fitted
$\staleage$ effect is \emph{not} interpretable as the true renewal
profile (it mixes the genuine gap dependence with the contamination).
The residual bias after all three mitigations is signed: fitted hazards
remain biased \emph{downward}, most severely at large staleness ages,
and every probability this chapter emits for a long-quiet firm is to be
read as a lower bound in that direction.

\begin{protocol}[Panel construction and audits]\label{prot:ch09:panel}
The panel builder emits, per horizon: the grid \cref{eq:ch09:grid}; rows
filtered by $\entry_i \le t_b < \Tend$, the exposure-window rule with
recorded $E$, and $\atrisk_i(t_b)=1$; labels \cref{eq:ch09:label} for
fully exposed buckets and censor flags with $e_{i,b}$ otherwise;
features exactly as \cref{def:ch09:panel} under \cref{prot:ch02:pit},
with the exp26 prohibition on end-of-sample snapshot columns (the
companies-table \code{last\_*} fields, totals, employees) enforced by
schema. Audits, run on every build: (a) no row precedes its firm's
entry; (b) no label consults data past $\Tend$; (c) censor-flag
retention --- the count of censored rows must equal the count implied by
the grid and $\Tend$; (d) as-of-join look-ahead test on macro features;
(e) cross-horizon coherence --- empirical event rates must be
nondecreasing in $\horizon$ on shared origins; (f) base rate and row
count by era, published next to every fit.
\end{protocol}

\section{Dependence: what the rows share}\label{ch09:sec:dependence}

The panel looks like a supervised-learning table and is not an i.i.d.\
sample. Rows share firms across buckets: the same latent quality that
\cref{ch:stale} models as a persistent state makes $y_{i,b}$ and
$y_{i,b'}$ positively dependent for all pairs of buckets of the same
firm. Rows share time within buckets: macro regimes and sector cycles
are common shocks, and the covariates-dominate verdict of
\cref{app:results} is precisely the statement that this shared
structure carries most of the signal. Both dependencies survive
conditioning on the features whenever the features are incomplete ---
which defect (C5) guarantees they are. The effective sample size is
therefore far smaller than the row count, nominal i.i.d.\ standard
errors are fictions, and shuffled cross-validation is leakage by
construction (rows from the same firm and era land on both sides of the
fold).

\begin{warning}\label{warn:ch09:sliding}
Sliding the bucket window (origins every $s < \horizon$ weeks) does not
create data. Each realized deal then appears in $\lceil \horizon/s
\rceil$ positive labels, and overlapping windows re-count the same
survival exposure: the resulting objective is not the likelihood of any
process, and any resampling scheme that treats the rows as exchangeable
understates uncertainty by a factor up to $\sqrt{\horizon/s}$. Sliding
windows are pseudo-replication, not augmentation.
\end{warning}

\begin{decision}{D9.3 --- disjoint buckets for training, sliding for
evaluation only}
Training panels use disjoint (non-overlapping) buckets per horizon, as
in \cref{eq:ch09:grid}; the likelihood factorization of
\cref{thm:ch09:hazardlink}(i) is exact only then. Phase-shifted disjoint
grids (origins offset by $s < \horizon$) may be built as robustness
replicates but are never pooled into one training set. Sliding origins
\emph{are} permitted for evaluation curves --- lift@$K$ or calibration
as a function of the decision date --- because evaluation is
description, not estimation; their inference uses the blocked resampling
below. Rationale: disjointness trades a constant factor of rows for an
exact likelihood and honest uncertainty; the rehearsal's era-blocked
evaluations showed the between-era variance dominates anyway, so the
discarded overlap bought little. Reversal trigger: if the $\horizon=104$
panel's bucket count per firm ($\approx 6$ over the 2011--2024 window)
proves too coarse for stable fitting, revisit pooled phase-shifted grids
with firm-and-era clustered inference priced in.
\end{decision}

Uncertainty for every panel-derived quantity --- metrics, calibration
curves, coefficient contrasts --- is computed by resampling that
respects both sharing directions: firms are resampled as whole clusters
(all rows of a firm move together) and eras as contiguous blocks (all
rows of a calendar block move together), following the temporal split
and blocked-inference protocol fixed in \cref{ch:gof}. The split itself
is temporal: training era, validation era (model selection, early
stopping, calibration), test era touched once. This is the same
discipline the exp26 benchmark harness already implements ---
temporal early stopping strictly inside the training era, a single
touch of the test era --- and the panel stack reuses that procedure
under the calendar-era boundaries fixed by \cref{ch:gof}.

\section{Class imbalance and negative sampling}\label{ch09:sec:sampling}

Positives are rare: most at-risk firm-buckets contain no deal, with the
base rate rising in $\horizon$. Rarity has two costs --- a statistical
one for the parametric benchmarks and a computational one for the full
panel --- and both have exact treatments; neither is an excuse for
folklore reweighting.

For the linear cloglog and logistic benchmarks that every boosted model
must beat, rare-event finite-sample bias is real: the MLE intercept is
biased in small-event-count regimes and separation becomes possible in
sparse eras. The corrections are standard and committed: the
\citet{kingzeng2001} rare-events adjustment for logistic fits, and
\citet{firth1993} penalization whenever any era-by-feature cell is
sparse enough to threaten separation.

For scale, the committed device is case-control subsampling of
negatives, with its exact correction.

\begin{proposition}[Subsampling offset]\label{prop:ch09:offset}
Retain every positive row and each negative row independently with
probability $r \in (0,1]$, with retention independent of $\zfeat$ given
$y$; let $S=1$ denote retention. If $p(\zfeat) = \Prob(y=1 \st \zfeat)$,
then
\begin{equation}\label{eq:ch09:subsample}
\operatorname{logit} \Prob\big(y = 1 \st \zfeat,\, S=1\big)
\;=\; \operatorname{logit} p(\zfeat) \;-\; \log r .
\end{equation}
Hence if the population model is logistic with predictor
$\alpha + f(\zfeat)$, the retained sample follows a logistic model with
the same $f$ and intercept $\alpha - \log r$: fitting on the subsample
with the fixed offset $-\log r$ in the predictor (equivalently, adding
$\log r$ to every fitted margin at prediction time) recovers the
population model exactly.
\end{proposition}

\begin{proof}
By Bayes,
$\Prob(y=1 \st \zfeat, S=1) = p(\zfeat) \cdot 1 \,\big/\,
\big[ p(\zfeat) \cdot 1 + (1-p(\zfeat))\, r \big]$,
so the retained odds are $p(\zfeat) / \big[(1-p(\zfeat))\, r\big]$, i.e.\
the population odds divided by $r$; taking logs gives
\cref{eq:ch09:subsample}. The second claim follows because the shift is
constant in $\zfeat$ and the logistic family is closed under intercept
shifts.
\end{proof}

Two remarks bound the result's reach. First, the exactness is a property
of the \emph{logistic} link: the retained-sample law under a cloglog
population model is not a cloglog model with shifted intercept, so
subsampling lives in the logistic pipeline, where
\cref{thm:ch09:hazardlink}(iii) already prices the link discrepancy, and
the calibration layer absorbs the residual. Second, the same logic
underwrites the ranking stack: estimating a conditional logit on a
\emph{sampled subset of alternatives} is consistent when the sampling
protocol satisfies the positive-conditioning property, by exactly the
same Bayes correction \citep{mcfadden1978} --- the sampled-softmax
training of \cref{ch:ranking} and \cref{eq:ch09:subsample} are one
result in two costumes: a known sampling scheme induces a computable
likelihood correction, so sampling is free of bias whenever it is
honestly recorded.

For the boosted trees themselves the folklore alternative is
\code{scale\_pos\_weight}: multiply positive-class gradients by a
constant $w$. We reject it for probability work. In the unregularized,
infinite-capacity limit weighting shifts every logit by $\log w$
uniformly --- harmless. In an actual boosted tree, leaf values are
$-G/(H+\lambda_{\mathrm{reg}})$ (\cref{ch09:sec:xgb}) and splits are
accepted against $\gamma$ and \code{min\_child\_weight} thresholds on
$H$: reweighting rescales $G$ and $H$ heterogeneously across leaves,
so the induced distortion is \emph{not} a constant shift, and no single
intercept correction undoes it. Subsampling with the offset of
\cref{prop:ch09:offset}, followed by post-hoc calibration on the
\emph{unsubsampled} validation era, achieves the same variance-per-cost
benefit with a distortion we can state and remove.

\begin{decision}{D9.4 --- subsample negatives, correct exactly,
calibrate after}
Training panels may subsample negatives at a recorded rate $r$ per
era-stratum; positives are never subsampled. Corrections are mandatory
and layered: the $\log r$ margin correction of \cref{prop:ch09:offset}
at prediction time, then the calibration map of
\cref{ch09:sec:calibration} fit on the full (unsubsampled) validation
era. \code{scale\_pos\_weight} and other loss reweightings are banned in
the probability pipeline (permitted in pure-ranking experiments, where
only order matters). Rationale: \cref{prop:ch09:offset} makes
subsampling exactly correctable; reweighting is not. Reversal trigger:
none; this is a mathematical, not empirical, commitment.
\end{decision}

\section{XGBoost as the committed learner}\label{ch09:sec:xgb}

The learner attached to the panel is gradient-boosted trees in the
XGBoost formulation \citep{chenguestrin2016}, i.e.\ the functional
gradient view of \citet{friedman2001} with a second-order objective and
explicit regularization. We derive the fitting rule once, because two of
its internals (the Hessian denominators and the default directions) are
load-bearing for our data.

The model after $m$ rounds is $F_m(\zfeat) = \sum_{k \le m}
f_k(\zfeat)$, each $f_k$ a regression tree with $T$ leaves, leaf weights
$w \in \R^T$, and leaf-assignment map $q$. Round $m$ solves
\begin{equation}\label{eq:ch09:xgbobj}
\min_{f} \;\; \sum_{n} l\big(y_n,\; F_{m-1}(\zfeat_n) + f(\zfeat_n)\big)
\;+\; \gamma T \;+\; \tfrac{1}{2}\lambda_{\mathrm{reg}} \norm{w}^2 ,
\end{equation}
with $l$ the per-row loss --- for us \code{binary:logistic},
$l(y, F) = -y F + \log(1 + e^{F})$, or the cloglog-exposure loss implied
by \cref{eq:ch09:offsetlik} as a custom objective.

\begin{proposition}[Second-order boosting step]\label{prop:ch09:xgb}
Let $g_n = \partial_F\, l(y_n, F)\big|_{F_{m-1}(\zfeat_n)}$ and
$h_n = \partial^2_F\, l(y_n, F)\big|_{F_{m-1}(\zfeat_n)}$, and for a
fixed tree structure $q$ let $I_j = \{n : q(\zfeat_n) = j\}$,
$G_j = \sum_{n \in I_j} g_n$, $H_j = \sum_{n \in I_j} h_n$. To second
order in $f$, the objective \cref{eq:ch09:xgbobj} is minimized over leaf
weights by
\[
w_j^{*} \;=\; -\,\frac{G_j}{H_j + \lambda_{\mathrm{reg}}},
\qquad
\text{with optimal value}
\quad
-\tfrac{1}{2} \sum_{j=1}^{T} \frac{G_j^{2}}{H_j + \lambda_{\mathrm{reg}}}
\;+\; \gamma T,
\]
and the gain of splitting a leaf with statistics $(G, H)$ into
$(G_L, H_L)$ and $(G_R, H_R)$ is
\begin{equation}\label{eq:ch09:gain}
\mathrm{Gain} \;=\; \tfrac{1}{2}\left[
\frac{G_L^{2}}{H_L + \lambda_{\mathrm{reg}}}
+ \frac{G_R^{2}}{H_R + \lambda_{\mathrm{reg}}}
- \frac{(G_L + G_R)^{2}}{H_L + H_R + \lambda_{\mathrm{reg}}}
\right] - \gamma .
\end{equation}
For the logistic loss, $g_n = \sigma(F_{m-1}(\zfeat_n)) - y_n$ and
$h_n = \sigma(F_{m-1}(\zfeat_n))\big(1 - \sigma(F_{m-1}(\zfeat_n))\big)$.
\end{proposition}

\begin{proof}
Taylor-expanding each loss term to second order,
$l(y_n, F_{m-1} + f) \approx l(y_n, F_{m-1}) + g_n f(\zfeat_n) +
\tfrac12 h_n f(\zfeat_n)^2$, and summing within leaves (where
$f \equiv w_j$) turns \cref{eq:ch09:xgbobj} into $\sum_j [\, G_j w_j +
\tfrac12 (H_j + \lambda_{\mathrm{reg}}) w_j^2 \,] + \gamma T$ up to a
constant: a separable quadratic, minimized at $w_j^{*}$ with the stated
value. The gain formula is the difference of optimal values before and
after the split, with $\gamma$ charged for the extra leaf
\citep{chenguestrin2016}.
\end{proof}

\paragraph{Why trees earn their keep here.} Four reasons, one of them a
measured verdict. First, \emph{the staleness interaction is the model}:
\cref{ch:stale} establishes that the predictive content of a LOCF
covariate decays with its age, so the truth has the form
$f(\tilde\xcov, \staleage)$ with effect modification in $\staleage$ ---
exactly what axis-aligned splits represent cheaply and what a linear
cloglog must hand-engineer term by term. Second, PitchBook columns are
incorrigibly mixed-type: round-type categories, flags, sizes spanning
orders of magnitude, and pervasive missingness; XGBoost's sparsity-aware
default directions \citep{chenguestrin2016} learn, per split, which way
missing values go --- and since missingness here is informative
(undisclosed size is a signal, not an accident), learning the direction
is modeling, not imputation. Third, nonlinearity in the engineered
history (trailing counts, raised-to-date) is expected and cheap to
capture. Fourth, and decisively: the B1 benchmark verdict of
\cref{app:results}. At strict feature parity --- the same six features,
the same choice sets, the same temporal split --- XGBoost beat the linear
conditional logit by $+0.22$ nats and $+9$pp top-5 on regime A. The
linear-utility restriction is thereby empirically falsified on the
identity target even in a clean six-feature world; on the panel, whose
feature set is wider and dirtier, we do not expect linearity to fare
better. The linear cloglog GLM remains in the harness as the benchmark
every boosted fit must beat, but it is the challenger, not the champion.

\paragraph{Hyperparameters that matter.} Depth (\code{max\_depth})
bounds interaction order; the rehearsal's gains appeared already at
depth 3--5, and deeper trees mostly buy overfitting to firm identity.
\code{min\_child\_weight} thresholds the leaf Hessian mass $H_j$; under
the logistic loss $h_n = p_n(1-p_n)$, so the threshold is an
effective-count guard that, under class imbalance, mechanically prevents
leaves fit to a handful of positives --- it is the imbalance safeguard
that \emph{is} compatible with calibration, unlike loss reweighting.
Row subsampling per boosting round reduces variance but should sample
firms as clusters when feasible, for the reasons of
\cref{ch09:sec:dependence}; column subsampling decorrelates trees in the
presence of the many collinear trailing-window counts. The number of
trees is chosen by early stopping on the temporal validation era ---
never by shuffled cross-validation --- following the exp26 protocol
(grid over depth and learning rate, early stop on the validation tail,
refit at the chosen round count, one shot at test).

\paragraph{Monotonicity constraints.} XGBoost accepts per-feature
monotone constraints, and economics suggests candidates: probability
nondecreasing in trailing deal counts, nonincreasing in the policy rate.
The tempting constraint --- intensity decreasing in staleness age
$\staleage$ --- we explicitly decline to impose. The true gap dependence
is not monotone: renewal-style maturation pushes hazard \emph{up} over
the first year-plus of a gap before frailty heterogeneity pulls the
population hazard down \citep{vaupel1979,aalen2008}, and the apparent
steep decline at large $\staleage$ is partly the (C3) contamination
artifact of \cref{ch09:sec:censoring}, not economics. Forcing
monotonicity would hard-code the artifact. Constraints are therefore
used as robustness probes --- fit with and without, compare test loss
and reliability --- and imposed only where the constrained fit is not
measurably worse.

\section{Probability calibration}\label{ch09:sec:calibration}

The decision layer of \cref{ch:decision} consumes the panel's outputs as
probabilities: expected shortlist yield $\sum_{i \in \text{list}}
\hat p_i$, thresholds set by expected utility, horizon-qualified
pipeline pacing. A model whose ranking is superb but whose probabilities
run hot by a factor of two poisons every one of those computations,
which is why calibration here is mandatory infrastructure, not a
finishing touch. Three distortions are anticipated and named: the
logit--cloglog gap of \cref{thm:ch09:hazardlink}(iii), residual
subsampling distortion after the $\log r$ correction of
\cref{prop:ch09:offset}, and the boosted-margin distortion documented by
\citet{niculescumizil2005}, who find boosted trees among the models that
benefit most from post-hoc mapping.

The committed instruments are the two standard ones. \emph{Platt
scaling} \citep{platt1999}: fit $c(m) = \sigma(Am + B)$ on validation
margins by logistic regression --- two parameters, stable on modest
validation eras, and well matched to distortions that are smooth and
sigmoidal in the margin. \emph{Isotonic regression}
\citep{zadroznyelkan2002}: the nonparametric monotone map fit by pooled
adjacent violators --- strictly more flexible, needs materially more
validation data, and prone to piecewise-constant artifacts in the tails
where our decisions live. Guidance follows
\citet{niculescumizil2005}: Platt as default; isotonic when the
validation reliability curve is visibly non-sigmoidal \emph{and} the
validation era supplies enough positives per bin to estimate it ---
per-horizon, the 26-week panel is the most likely to qualify, the
104-week panel the least.

\begin{protocol}[Calibration and its audit]\label{prot:ch09:calib}
(a) The calibration map is fit on the validation era only, after the
$\log r$ subsampling correction, and is part of the model artifact: no
downstream consumer ever sees raw margins. (b) Reliability diagrams
(binned $\hat p$ versus empirical frequency, with firm-cluster,
era-block resampled bands) are produced per horizon and per era, before
and after calibration. (c) Scores: the Brier score with its
reliability--resolution--uncertainty decomposition, and the log score;
both are proper, so no calibration map can game them in expectation
\citep{gneitingraftery2007}. (d) Drift: calibration is re-audited on
each new era; a reliability shift with stable discrimination triggers
recalibration, not refitting.
\end{protocol}

\begin{decision}{D9.5 --- calibration is mandatory}
Every deployed horizon model carries a calibration map fit per
\cref{prot:ch09:calib}; uncalibrated scores are internal quantities.
Rationale: the consuming actions are probability arithmetic
(\cref{ch:decision}), and the three distortion sources above are known
and repairable at the cost of a validation-era fit. Reversal trigger:
none. If a future learner is verified reliably calibrated end-to-end,
the map degenerates to the identity, at zero cost.
\end{decision}

\section{Survival-aware boosting alternatives}\label{ch09:sec:aft}

The per-$\horizon$ classifier design pays two visible costs: censored
buckets are dropped (\cref{ch09:sec:censoring}), losing the most recent
eras precisely at the long horizons; and the three horizon models are
fit separately, so nothing enforces coherence of
$\hat p_i(t, 26) \le \hat p_i(t, 52) \le \hat p_i(t, 104)$ beyond the
audit in \cref{prot:ch09:panel}. Both costs are removed by boosting a
survival model instead of a classifier.

\emph{XGBoost--AFT} \citep{barnwal2022} boosts an accelerated failure
time model, $\log W_i = \mu(\zfeat_i) + \sigma \eps$, with $\eps$
standard normal, logistic, or extreme-value, and $\mu$ the boosted
ensemble. Its likelihood consumes exactly the data the classifier
discards: an observed waiting time contributes a density term, a bucket
that crosses $\Tend$ contributes a right-censored survival term
$\Prob(W > \Tend - t_b)$, and grouped observations contribute
interval-censored terms $\Prob(l < W \le u)$ --- censored rows are
retained natively, no offsets required. One fitted model produces the
whole curve $\horizon \mapsto \hat\Prob(W \le \horizon)$, monotone in
$\horizon$ by construction, serving all three horizons and any future
one. \emph{Gradient-boosted Cox} \citep{polsterl2020} boosts the partial
likelihood instead: it inherits the proportional-hazards reading of
\cref{thm:ch09:hazardlink}(i) and needs no error-distribution choice,
but it delivers relative risks; absolute horizon probabilities require a
Breslow-type baseline estimate on top, at which point its calibration
story is no simpler than the classifier's.

The costs run the other way. AFT buys its coherent curve with a global
parametric assumption on $\eps$: a misspecified error shape miscalibrates
\emph{every} horizon in a correlated way, and repairing it per horizon
reintroduces exactly the per-$\horizon$ calibration maps the design was
meant to avoid. AFT also consumes exact waiting times, thereby
re-inheriting the date-jitter component of (C4) that the bucket design
is immune to by construction. And the surrounding tooling ---
imbalance-aware sampling theory (\cref{prop:ch09:offset} is
logistic-specific), calibration instruments, evaluation conventions ---
is mature for binary classification and improvised for boosted AFT.

\begin{decision}{D9.6 --- per-$\horizon$ classifiers now, AFT as the
benchmarked alternative}
The committed T3 stack is the per-$\horizon$, cloglog-consistent
classifier pipeline of this chapter. XGBoost--AFT is maintained in the
benchmark harness on the identical feature set and splits, evaluated on
the same per-horizon probabilities $\hat\Prob(W \le \horizon)$ after its
own calibration audit. Rationale: exact censoring and sampling theory,
mature calibration, and (C4) robustness today outweigh the coherent
curve and censored-row retention. Reversal trigger: if AFT beats the
calibrated per-$\horizon$ classifiers on test log score at two of three
horizons in two consecutive evaluation campaigns --- or if D9.2's
reversal fires and the decision layer demands curves --- AFT is
promoted.
\end{decision}

The repository's prior art is acknowledged and mined: experiment
\code{exp16} already implements a discrete-time logistic-hazard model
with Kaplan--Meier calibration reference and a pairwise ranking-loss
hybrid on synthetic firms --- the panel stack of this chapter is that
sketch industrialized, with the labeling, censoring, and sampling theory
made explicit.

\section{Evaluation, and the bridge to ranking}\label{ch09:sec:eval}

Evaluation follows the temporal split of \cref{ch:gof}, reuses the exp26
harness conventions (selection on the validation era, one shot at test),
and reports per horizon and per era. The headline metrics are chosen to
proxy the consuming action, per \cref{eq:ch01:voi}. \emph{PR-AUC}, not
ROC-AUC, is the discrimination summary: under heavy imbalance ROC-AUC is
inflated by the ocean of easy negatives, while precision--recall tracks
the operating region the sourcing action lives in; because PR-AUC
depends on the base rate, it is reported next to the era's base rate
and never compared across eras or horizons without that context.
\emph{Lift at $K$} --- the event rate in the top-$K$ list divided by the
base rate, at $K$ set to realistic diligence capacities --- is the
closest offline proxy to sourcing value. \emph{Reliability and proper
scores} per \cref{prot:ch09:calib} complete the set: discrimination
without calibration is half a model here. All uncertainty is
firm-clustered, era-blocked (\cref{ch09:sec:dependence}).

Sorting firms by $\hat p_i(t, \horizon)$ is a ranking, and the T1
machinery of \cref{ch:ranking} is close enough that the relationship
must be stated rather than left to intuition. The two targets differ in
conditioning: the ranker estimates $\Prob(i \st \text{deal in } s, t,
\Hist)$ --- identity given that a deal happens, timing canceled ---
while the horizon model estimates a marginal that mixes identity with
timing mass; under the two-layer factorization of
\cref{ch:question}, heuristically $p_i(t,\horizon) \approx 1 -
\exp(-\int_t^{t+\horizon} \Lam_s(u)\, \pi_i(u) \dd u)$, so a
mid-ranked firm in a hot sector can out-probability a top-ranked firm
in a cold one. They differ in universe: the ranker competes firms
within a sector risk set and must own the outside option (C6); the
horizon model scores each tracked firm absolutely, no competition and
no outside option. And they differ in unit: recall@$K$ is computed
against realized deals and is calibration-free; lift@$K$ on the panel
is computed on all at-risk rows and consumes calibrated levels. The two
are complements, not rivals: the sector-conditional ranker is the
defensible object when timing is untrusted (\cref{ch:question}), the
firm-level horizon model is the object the diligence trigger consumes,
and the composed continuous-time stack furnishes a third estimate of
the same $p_i(t,\horizon)$ by integrating the fitted ground intensity
against the mark factor. Comparing the boosted panel's
$\hat p_i(t,\horizon)$ with that composed structural estimate on
identical origins is the standing cross-stack audit demanded by design
decision D1.2: when the two stacks disagree systematically, one of them
is wrong about something nameable --- and every diagnostic in this
chapter exists to name it.

\begin{codehooks}
\emph{Panel layer} (to be written): \modrepo{panel/builder.py} ---
grids, at-risk filtering, labels, exposure fractions $e_{i,b}$, censor
flags, and feature assembly per \cref{def:ch09:panel} under
\cref{prot:ch02:pit}; \modrepo{panel/audits.py} --- the audit battery of
\cref{prot:ch09:panel} (entry consistency, censor-flag retention,
look-ahead tests, cross-horizon monotonicity, base rates by era).
\emph{Model layer} (to be written): \modrepo{models/xgb\_horizon.py} ---
per-$\horizon$ XGBoost with the D9.4 subsampling-offset correction, an
optional cloglog-exposure custom objective implementing
\cref{eq:ch09:offsetlik}, monotone-constraint configuration, and the
calibration map of \cref{prot:ch09:calib} as part of the persisted
artifact; \modrepo{models/xgb\_aft.py} --- the D9.6 AFT benchmark on
identical features and splits. \emph{Harness}: the temporal-split,
validation-tail early-stopping, grid-then-one-test-shot protocol is
inherited from \code{experiments/exp26\_xgb\_benchmark.py} and must not
be reimplemented; \code{experiments/exp16\_survival.py} (discrete-time
logistic hazard, KM calibration, ranking-loss hybrid) is the prior art
this chapter supersedes. Linear benchmarks (cloglog GLM with
\citet{firth1993} penalization) live beside the boosted models in the
same harness so that every gain is measured, as in the fired B1 test,
at feature parity.
\end{codehooks}
